\chapter{NetworkModel come adapter UI}
\label{ch:network-adapter}

\section{Obiettivo (Fase 7)}

Separare la \textbf{shell osservabile} usata dall'interfaccia dalla logica
di query sulla topologia. \texttt{NetworkModel} resta nel target app;
le operazioni pure vivono in \module{FDCDomain}.

\section{NetworkTopologyService}

\texttt{NetworkTopologyService} (\repo{Domain/Services/}) è una struct
\texttt{Sendable} che incapsula \texttt{RailwayTopology} e offre:

\begin{itemize}
  \item lookup: \texttt{findNode}, \texttt{findEdge}
  \item filtri: \texttt{allStations}, \texttt{allJunctions}
  \item ordinamento: \texttt{sortedNodes}, \texttt{sortedEdges}
  \item connettività: \texttt{getConnectedNodeIds}, \texttt{getNeighborStations}
  \item geometria: \texttt{calculateDistance} (hub-aware, km)
\end{itemize}

Nessuna mutazione, nessuna dipendenza da SwiftUI o undo.

\section{NetworkModel}

\texttt{NetworkModel} delega le query a \texttt{NetworkTopologyService}
tramite una computed property \texttt{topologyService}. Espone
\texttt{topology: RailwayTopology} per i servizi di scheduling.

Le mutazioni (aggiunta nodi, binari, undo) restano in
\repo{Models/Network/NetworkModel+Updates.swift}.

\section{DomainBridge}

\texttt{DomainBridge.swift} espone il package e risolve l'ambiguità
\texttt{Edge} (SwiftUI vs dominio):

\begin{verbatim}
@_exported import FDCDomain
public typealias Edge = FDCDomain.Edge
\end{verbatim}

I typealias legacy per gli altri tipi sono stati rimossi.

\section{TrainCategory e DTO}

\begin{itemize}
  \item \texttt{TrainCategory} --- enum Foundation in \repo{Domain/Model/};
        estensione SwiftUI in \repo{Models/TrainCategory+UI.swift}
  \item \texttt{RailwayNetworkDTO} --- struct \texttt{Codable} in
        \repo{Infrastructure/} (I/O JSON, non dominio puro)
\end{itemize}

\section{Test}

\begin{itemize}
  \item \textbf{2 test} aggiuntivi in \texttt{PathfindingTests} (app)
  \item \textbf{2 test} in \texttt{FDCDomainTests} per \texttt{NetworkTopologyService}
\end{itemize}

\section{Prossimi passi (Fase 8)}

\begin{itemize}
  \item Estrarre \texttt{FDCScheduling} in package SPM
  \item Protocolli in \texttt{Domain/Services/} per ConflictManager, GA, AI
\end{itemize}
